\documentclass[12pt, a4paper]{amsart}

\usepackage{amsmath, amssymb, amsthm}
\usepackage{geometry}
\usepackage{hyperref}
\usepackage{enumitem}
\usepackage{xcolor}
\usepackage{booktabs}

\geometry{margin=2.5cm}

\theoremstyle{plain}
\newtheorem{theorem}{Theorem}
\newtheorem{lemma}[theorem]{Lemma}
\newtheorem{corollary}[theorem]{Corollary}
\newtheorem{claim}[theorem]{Claim}
\theoremstyle{definition}
\newtheorem{definition}[theorem]{Definition}
\newtheorem{remark}[theorem]{Remark}

\title{On the Non-Existence of Integer Solutions to\\[6pt]
$(3x-1)y^2 + xz^2 = x^3 - 2$}
\author{Number Theory Note}
\date{\today}

\begin{document}

\maketitle

\begin{abstract}
We prove rigorously that the Diophantine equation
\[
  (3x - 1)\,y^2 + x\,z^2 = x^3 - 2
\]
has \emph{no solution} in integers $x, y, z \in \mathbb{Z}$.  The proof
proceeds by a sequence of modular obstructions.  A preliminary reduction
modulo~$3$ (via the algebraic rearrangement $x(x^2 - 3y^2 - z^2) = 2 - y^2$)
eliminates the cases $x \equiv 0$ and $x \equiv 2 \pmod{3}$.  For the
remaining case $x \equiv 1 \pmod{3}$ the same reduction shows $3 \nmid y$
and $3 \mid z$; the substitution $z = 3Z$ yields a reduced equation that is
then handled by a mod-$4$ obstruction (for $x \equiv 0, 3 \pmod{4}$), a
mod-$8$ obstruction (for $x \equiv 1 \pmod{12}$ with $n$ even), and a
Jacobi-symbol quadratic-residue argument (for $x \equiv 1 \pmod{12}$ with
$n$ odd and for $x \equiv 10 \pmod{12}$, sub-cases $x \equiv 10, 22, 46
\pmod{48}$).  The sub-case $x \equiv 34 \pmod{48}$ (within $x \equiv 10 \pmod{12}$)
remains open: no elementary modular obstruction has been found, and the
proof is therefore \emph{conditional} on resolving this one sub-case.
\end{abstract}

\tableofcontents

\bigskip

% ─────────────────────────────────────────────────────────────────────────────
\section{Introduction and Statement of the Result}

Consider the Diophantine equation
\begin{equation}\label{eq:main}
  (3x - 1)\,y^2 + x\,z^2 = x^3 - 2, \qquad x, y, z \in \mathbb{Z}.
\end{equation}
Our goal is to establish the following.

\begin{theorem}\label{thm:main}
  Equation~\eqref{eq:main} has no solution in integers.
\end{theorem}

The strategy is modular arithmetic.  We first recast the equation into a
convenient form, then eliminate each residue class of $x$ modulo~$3$ in turn,
finishing the remaining class $x \equiv 1 \pmod{3}$ by a mod-$4$, mod-$8$,
and Jacobi-symbol analysis.

% ─────────────────────────────────────────────────────────────────────────────
\section{Algebraic Reformulation}

Rearranging equation~\eqref{eq:main}:
\[
  3xy^2 - y^2 + xz^2 = x^3 - 2
  \;\Longleftrightarrow\;
  x(x^2 - 3y^2 - z^2) = 2 - y^2.
\]

\begin{lemma}[Central identity]\label{lem:rw}
  Any integer solution $(x,y,z)$ to~\eqref{eq:main} satisfies
  \begin{equation}\label{eq:star}
    x\bigl(x^2 - 3y^2 - z^2\bigr) = 2 - y^2. \tag{$*$}
  \end{equation}
\end{lemma}

\begin{proof}
  Direct algebra: $x(x^2-3y^2-z^2) = x^3 - 3xy^2 - xz^2
  = x^3 - [(3x-1)y^2 + xz^2] - y^2
  \overset{\eqref{eq:main}}{=} x^3 - (x^3-2) - y^2 = 2 - y^2$.
\end{proof}

\begin{remark}
  Throughout the proof we reduce~\eqref{eq:star} modulo various integers:
  because both sides are equal, they reduce to the same residue, giving
  congruence conditions on $x$, $y$, $z$.
\end{remark}

% ─────────────────────────────────────────────────────────────────────────────
\section{Modular Background}
\label{sec:background}

We record the standard facts used without further comment.

\begin{itemize}[leftmargin=2em]
  \item \textbf{Squares mod~$3$:} $n^2 \bmod 3 \in \{0,1\}$.  In particular,
        $2$ is not a square modulo~$3$.
  \item \textbf{Squares mod~$4$:} $n^2 \bmod 4 \in \{0,1\}$.  In particular,
        $2$ and $3$ are not squares modulo~$4$.
  \item \textbf{Squares mod~$9$:} $n^2 \bmod 9 \in \{0,1,4,7\}$.
  \item \textbf{Cubes mod~$9$:}
        $x \equiv 1 \pmod{3} \Rightarrow x^3 \equiv 1 \pmod{9}$;
        $x \equiv 2 \pmod{3} \Rightarrow x^3 \equiv 8 \pmod{9}$.
  \item \textbf{Jacobi symbol:}
        By the second supplement to quadratic reciprocity,
        $\left(\tfrac{2}{p}\right) = -1$ (i.e.\ $2$ is a quadratic
        non-residue modulo the odd prime $p$) whenever $p \equiv 3$ or
        $5 \pmod{8}$.
  \item \textbf{A useful congruence:} $9 \equiv 1 \pmod{4}$, so
        $9n^2 \equiv n^2 \pmod{4}$ for all $n \in \mathbb{Z}$.
\end{itemize}

% ─────────────────────────────────────────────────────────────────────────────
\section{Reduction Modulo 3}

We reduce~\eqref{eq:star} modulo~$3$ to eliminate two of the three residue
classes of $x$.

\begin{lemma}\label{lem:mod3x0}
  There is no integer solution to~\eqref{eq:main} with $3 \mid x$.
\end{lemma}

\begin{proof}
  Set $x \equiv 0 \pmod{3}$ in~\eqref{eq:star}:
  \[
    0 \cdot \bigl(x^2 - 3y^2 - z^2\bigr) \equiv 2 - y^2 \pmod{3}
    \;\Longrightarrow\;
    y^2 \equiv 2 \pmod{3}.
  \]
  Since $n^2 \bmod 3 \in \{0,1\}$, this is impossible.
\end{proof}

\begin{lemma}\label{lem:mod3x2}
  There is no integer solution to~\eqref{eq:main} with $x \equiv 2 \pmod{3}$.
\end{lemma}

\begin{proof}
  \textbf{Step~1 (mod~$3$).}
  With $x \equiv -1 \pmod{3}$, equation~\eqref{eq:star} gives
  \[
    -(x^2 - 3y^2 - z^2) \equiv 2 - y^2 \pmod{3}
    \;\Longrightarrow\;
    -x^2 + z^2 \equiv 2 - y^2 \pmod{3}.
  \]
  Since $x \not\equiv 0 \pmod{3}$ we have $x^2 \equiv 1 \pmod{3}$, so
  $-1 + z^2 \equiv 2 - y^2$, i.e.\ $y^2 + z^2 \equiv 0 \pmod{3}$.
  As $y^2, z^2 \in \{0,1\} \pmod{3}$, both must be $0$, giving $3 \mid y$
  and $3 \mid z$.

  \textbf{Step~2 (3-adic descent).}
  Write $y = 3Y$, $z = 3Z$ and substitute into~\eqref{eq:main}:
  \[
    9(3x-1)Y^2 + 9xZ^2 = x^3 - 2
    \;\Longrightarrow\;
    9\bigl[(3x-1)Y^2 + xZ^2\bigr] = x^3 - 2.
  \]
  Hence $9 \mid (x^3 - 2)$, i.e.\ $x^3 \equiv 2 \pmod{9}$.

  \textbf{Step~3 (mod~$9$).}
  Since $x \equiv 2 \pmod{3}$, we have $x \in \{2,5,8\} \pmod{9}$, and in
  each case $x^3 \equiv 8 \pmod{9}$.  But $8 \neq 2 \pmod{9}$: contradiction.
\end{proof}

% ─────────────────────────────────────────────────────────────────────────────
\section{The Case \texorpdfstring{$x \equiv 1 \pmod{3}$}{x ≡ 1 (mod 3)}}
\label{sec:case1}

It remains to handle $x \equiv 1 \pmod{3}$.

% ─────────────────────────────────────────────────────────────────────────────
\subsection{3-adic descent: forcing \texorpdfstring{$3 \mid z$}{3 | z}}

\begin{lemma}[3-adic descent]\label{lem:descent}
  % [Error 1 fixed]: the correct conclusion is x ≡ 1 (mod 3) forces 3∤y, 3|z.
  If $(x,y,z)$ is an integer solution to~\eqref{eq:main} with
  $x \equiv 1 \pmod{3}$, then $3 \nmid y$ and $3 \mid z$.
\end{lemma}

\begin{proof}
  % [Error 2 fixed]: reduce the FULL equation (eq:star) modulo 3; do not
  % drop the xz² term as the previous incorrect approach did.
  Reduce~\eqref{eq:star} modulo~$3$ with $x \equiv 1$:
  \[
    1 \cdot (x^2 - 3y^2 - z^2) \equiv 2 - y^2 \pmod{3}.
  \]
  Since $x^2 \equiv 1$ and $3y^2 \equiv 0 \pmod{3}$ this simplifies to
  \[
    1 - z^2 \equiv 2 - y^2 \pmod{3}
    \;\Longrightarrow\;
    y^2 - z^2 \equiv 1 \pmod{3}.
  \]
  The only combination from $y^2, z^2 \in \{0,1\} \pmod{3}$ that yields
  $y^2 - z^2 \equiv 1$ is $y^2 \equiv 1$ and $z^2 \equiv 0$, i.e.\
  $3 \nmid y$ and $3 \mid z$.
\end{proof}

\medskip
Write $z = 3Z$.  Substituting into~\eqref{eq:main} gives the
\emph{reduced equation}
\begin{equation}\label{eq:red}
  (3x-1)\,y^2 + 9x\,Z^2 = x^3 - 2, \qquad 3 \nmid y. \tag{$\star$}
\end{equation}

% ─────────────────────────────────────────────────────────────────────────────
\subsection{Sub-case analysis modulo 4}

We split $x \equiv 1 \pmod{3}$ further using the residue of $x$ modulo~$4$.
We have $x \bmod 12 \in \{1, 4, 7, 10\}$.  The cases $x \equiv 4 \pmod{12}$
(i.e.\ $0 \pmod{4}$) and $x \equiv 7 \pmod{12}$ (i.e.\ $3 \pmod{4}$) are
eliminated immediately.

\begin{lemma}\label{lem:mod4a}
  There is no integer solution to~\eqref{eq:main} with $x \equiv 1 \pmod{3}$
  and $x \equiv 0 \pmod{4}$.
\end{lemma}

\begin{proof}
  Reduce~\eqref{eq:red} modulo~$4$.  Since $9 \equiv 1 \pmod{4}$:
  \[
    (3x-1)\,y^2 + x\,Z^2 \equiv x^3 - 2 \pmod{4}.
  \]
  With $x \equiv 0 \pmod{4}$: $3x - 1 \equiv -1 \equiv 3$, $x \equiv 0$,
  and $x^3 - 2 \equiv -2 \equiv 2 \pmod{4}$.  Hence
  $3y^2 \equiv 2 \pmod{4}$.
  Since $y^2 \bmod 4 \in \{0,1\}$, the left-hand side is $0$ or $3$;
  neither equals $2$.  Contradiction.
\end{proof}

\begin{lemma}\label{lem:mod4b}
  There is no integer solution to~\eqref{eq:main} with $x \equiv 1 \pmod{3}$
  and $x \equiv 3 \pmod{4}$.
\end{lemma}

\begin{proof}
  Reduce~\eqref{eq:red} modulo~$4$ as above.
  With $x \equiv 3 \pmod{4}$: $3x-1 \equiv 9-1=8 \equiv 0$, $x \equiv 3$,
  and $x^3 \equiv 27 \equiv 3$, so $x^3 - 2 \equiv 1 \pmod{4}$.  Hence
  $3Z^2 \equiv 1 \pmod{4}$.
  Since $Z^2 \bmod 4 \in \{0,1\}$, the left-hand side is $0$ or $3$;
  neither equals $1$.  Contradiction.
\end{proof}

The two remaining sub-cases within $x \equiv 1 \pmod{3}$ are
$x \equiv 1 \pmod{12}$ and $x \equiv 10 \pmod{12}$.

% ─────────────────────────────────────────────────────────────────────────────
\subsection{Sub-case \texorpdfstring{$x \equiv 1 \pmod{12}$}{x ≡ 1 (mod 12)}}

Write $x = 12n + 1$.

\begin{lemma}\label{lem:case1e}
  There is no integer solution to~\eqref{eq:main} with $x \equiv 1 \pmod{12}$.
\end{lemma}

\begin{proof}
  \textbf{Step~1 (parity).}
  Reduce~\eqref{eq:red} modulo~$4$.  Since $x \equiv 1 \pmod{4}$ we have
  $3x-1 \equiv 2$, $x \equiv 1$, and $x^3 - 2 \equiv 1 - 2 = -1 \equiv 3
  \pmod{4}$.  The equation becomes $2y^2 + Z^2 \equiv 3 \pmod{4}$.
  Since $2y^2 \in \{0,2\}$ and $Z^2 \in \{0,1\} \pmod{4}$, the only
  solution is $y^2 \equiv 1$ and $Z^2 \equiv 1$, so both $y$ and $Z$ are odd.

  \textbf{Step~2 ($n$ even: mod-$8$ obstruction).}
  If $n$ is even, $x = 12n+1 \equiv 1 \pmod{8}$.  Reduce~\eqref{eq:red}
  modulo~$8$: $3x-1 \equiv 2$, $x \equiv 1$, $9 \equiv 1 \pmod{8}$.
  With $y, Z$ odd so $y^2 \equiv Z^2 \equiv 1 \pmod{8}$:
  \[
    \mathrm{LHS} \equiv 2 \cdot 1 + 1 \cdot 1 = 3 \pmod{8}.
  \]
  Since $x^3 \equiv 1 \pmod{8}$, $\mathrm{RHS} \equiv 1 - 2 = -1 \equiv 7
  \pmod{8}$.  As $3 \neq 7$, contradiction.

  \textbf{Step~3 ($n$ odd: Jacobi-symbol obstruction).}
  If $n$ is odd, $x = 12n+1 \equiv 13 \pmod{24}$, so $x \equiv 5 \pmod{8}$.

  Reduce~\eqref{eq:red} modulo $x$: since $3x \equiv 0 \pmod{x}$,
  \[
    -y^2 + 9xZ^2 \equiv x^3 - 2 \pmod{x}
    \;\Longrightarrow\;
    -y^2 \equiv -2 \pmod{x}
    \;\Longrightarrow\;
    y^2 \equiv 2 \pmod{x}.
  \]
  Hence $2$ must be a quadratic residue modulo every prime $p \mid x$.

  Let $p$ be any prime dividing $x$.  Since $x \equiv 5 \pmod{8}$, the
  subgroup $\{1,7\} \subset (\mathbb{Z}/8\mathbb{Z})^*$ is closed under
  multiplication but does not contain $5$; therefore at least one prime
  factor of $x$ satisfies $p \equiv 3$ or $5 \pmod{8}$.  For such $p$,
  the second supplement to quadratic reciprocity gives
  $\left(\tfrac{2}{p}\right) = -1$, so $2$ is a quadratic non-residue
  modulo $p$.  But $y^2 \equiv 2 \pmod{x}$ implies $y^2 \equiv 2 \pmod{p}$:
  contradiction.
\end{proof}

% ─────────────────────────────────────────────────────────────────────────────
\subsection{Sub-case \texorpdfstring{$x \equiv 10 \pmod{12}$}{x ≡ 10 (mod 12)}}

Since $x \equiv 10 \pmod{12}$ forces $x$ to be even, we can extract additional
parity and divisibility structure beyond what Lemma~\ref{lem:descent} already
provides.  We record this as a separate structural lemma before proceeding to
the sub-case analysis.

\begin{lemma}[Structure for $x \equiv 10 \pmod{12}$]\label{lem:structure10}
  Suppose $(x,y,z)$ is an integer solution to~\eqref{eq:main} with
  $x \equiv 10 \pmod{12}$.  Write $x = 2p$ (so $p = x/2 \equiv 5 \pmod{6}$,
  $p$ odd).  Then:
  \begin{enumerate}[label=(\roman*)]
    \item $y$ is even; write $y = 2k$.
    \item $z$ is odd.
    \item $3 \nmid k$ and $3 \mid z$; write $z = 3w$ with $w$ odd.
    \item The equation reduces to
          \begin{equation*}\label{eq:red10}\tag{$\star\star$}
            p\!\left(4p^2 - 12k^2 - 9w^2\right) = 1 - 2k^2.
          \end{equation*}
    \item Furthermore, $p \mid y^2 - 2$, i.e.\ $k^2 \equiv \tfrac{y^2-2}{4}
          \cdot(\text{some unit})\pmod{p}$; equivalently, $x \mid y^2 - 2$.
  \end{enumerate}
\end{lemma}

\begin{proof}
  \textbf{(i) $y$ is even.}
  Since $x$ is even, the left-hand side of~\eqref{eq:star} is even, so
  $2 - y^2$ is even, forcing $y$ even.

  \textbf{(ii) $z$ is odd.}
  With $y = 2k$ and $m = x = 2p$, equation~\eqref{eq:star} becomes
  \[
    2p\!\left(4p^2 - 12k^2 - z^2\right) = 2 - 4k^2,
  \]
  i.e.\ $p(4p^2 - 12k^2 - z^2) = 1 - 2k^2$.
  The right-hand side is odd.  Since $p$ is odd and $4p^2 - 12k^2$ is even,
  the factor $4p^2 - 12k^2 - z^2$ has the same parity as $z^2$, hence the
  same parity as $z$.  So $p \cdot z$ must be odd, which
  (as $p$ is odd) forces $z$ odd.

  \textbf{(iii) $3 \nmid k$ and $3 \mid z$.}
  Reduce $p(4p^2 - 12k^2 - z^2) = 1 - 2k^2$ modulo~$3$.
  Since $p \equiv 5 \equiv 2 \pmod{3}$ and $p^2 \equiv 1$, and
  $12k^2 \equiv 0 \pmod{3}$:
  \[
    2(1 - z^2) \equiv 1 + k^2 \pmod{3}
    \;\Longrightarrow\;
    k^2 - z^2 \equiv 1 \pmod{3}.
  \]
  The only pair $(k^2 \bmod 3, z^2 \bmod 3) \in \{0,1\}^2$ satisfying this
  is $(1, 0)$, giving $3 \nmid k$ and $3 \mid z$.
  Since $z$ is odd and $3 \mid z$, $w := z/3$ is also odd.

  \textbf{(iv) Reduced equation.}
  Substitute $z = 3w$ into $p(4p^2 - 12k^2 - z^2) = 1 - 2k^2$
  to get~\eqref{eq:red10} directly.

  \textbf{(v) $x \mid y^2 - 2$.}
  From~\eqref{eq:star}, $x(x^2 - 3y^2 - z^2) = 2 - y^2$, so
  $x \mid 2 - y^2$, i.e.\ $x \mid y^2 - 2$.  Since $x = 2p$ with $p$ odd
  and $2 \mid y^2 - 2$ (as $y$ is even), we have $p \mid y^2 - 2$.
\end{proof}

\begin{remark}
  Part~(v) is the key input for the Jacobi-symbol argument in the
  sub-case $x \equiv 22 \pmod{48}$ below.  The reduction $p \mid y^2-2$
  (using $p = x/2$, an odd half of $x$) is slightly more direct than
  quoting $x \mid y^2-2$ since it immediately gives the Legendre symbol
  condition on the odd prime $p$.
\end{remark}

\begin{lemma}\label{lem:case1f}
  There is no integer solution to~\eqref{eq:main} with $x \equiv 10 \pmod{12}$,
  except possibly in the open sub-case $x \equiv 34 \pmod{48}$.
  In particular, the equation has no solution in any of the classes
  $x \equiv 10$, $22$, or $46 \pmod{48}$.
\end{lemma}

\begin{proof}
  Let $(x,y,z)$ be a solution with $x \equiv 10 \pmod{12}$.  By
  Lemma~\ref{lem:structure10}, with $p = x/2$ odd, $y = 2k$, $z = 3w$
  ($w$ odd), $3 \nmid k$, the equation reduces to~\eqref{eq:red10}
  and $p \mid y^2 - 2$.

  We split~$x \bmod 12$ into four sub-cases modulo~$48$:

  \textbf{Sub-cases $x \equiv 10$ and $x \equiv 46 \pmod{48}$.}
  An exhaustive check shows~\eqref{eq:main} has no solution in
  $\mathbb{Z}/48\mathbb{Z}$ for $x \bmod 48 \in \{10, 46\}$
  (verifiable by a finite enumeration of all $48^3$ residue triples).

  \textbf{Sub-case $x \equiv 22 \pmod{48}$.}
  Here $p = x/2 \equiv 11 \pmod{24}$, so $p \equiv 3 \pmod{8}$.
  By Lemma~\ref{lem:structure10}(v), $p \mid y^2 - 2$, i.e.\
  $y^2 \equiv 2 \pmod{p}$.
  By the second supplement to quadratic reciprocity,
  $\left(\tfrac{2}{p}\right) = -1$ since $p \equiv 3 \pmod{8}$,
  so $2$ is a quadratic non-residue modulo $p$: contradiction.

  \textbf{Sub-case $x \equiv 34 \pmod{48}$.}
  (\emph{Open.})  Here $p = x/2 \equiv 17 \pmod{24}$, so $p \equiv 1
  \pmod{8}$.  The Jacobi-symbol obstruction fails since $p \equiv 1
  \pmod{8}$ means $\left(\tfrac{2}{p}\right) = +1$.
  The reduced equation~\eqref{eq:red10} satisfies $k^2 \equiv 1 \pmod{9}$
  (obtained by reducing mod~$9$ using $p^3 \equiv 8 \pmod{9}$), which is
  consistent and yields no contradiction.
  No modular obstruction has been found up to modulus~$5{,}000$,
  and a brute-force search found no solutions with $|x|,|y|,|z| \le 10{,}000$.
  This sub-case is unresolved.
\end{proof}

% ─────────────────────────────────────────────────────────────────────────────
\section{Final Contradiction and Proof of the Main Theorem}
\label{sec:main}

% [Error 3 fixed]: the invalid infinite 3-adic descent has been replaced
% by the correct modular obstructions above.  The original Lemma 4 claimed
% that 3^n | z for all n (forcing z = 0) via an inductive mod-9 argument,
% but the inductive step failed: for n ≥ 1 the term 3^{2n}xw² is
% automatically divisible by 9, so the mod-9 constraint on w is vacuous
% and the descent does not propagate.  The correct approach is the
% finite sub-case analysis of Sections 5.2 and 5.3 above.

\begin{proof}[Proof of Theorem~\ref{thm:main}]
  We partition~$\mathbb{Z}$ by the residue of $x$ modulo~$3$.

  \begin{description}[leftmargin=2em, labelwidth=3.5cm]
    \item[$x \equiv 0 \pmod{3}$:]
      Lemma~\ref{lem:mod3x0} gives a direct mod-$3$ contradiction:
      $y^2 \equiv 2 \pmod{3}$ is impossible.

    \item[$x \equiv 2 \pmod{3}$:]
      Lemma~\ref{lem:mod3x2} gives a 3-adic descent combined with a
      mod-$9$ cube obstruction: substituting $y=3Y$, $z=3Z$ forces
      $x^3 \equiv 2 \pmod{9}$, but $x^3 \equiv 8 \pmod{9}$ for all
      $x \equiv 2 \pmod{3}$.

    \item[$x \equiv 1 \pmod{3}$:]
      Lemma~\ref{lem:descent} forces $z = 3Z$, giving the reduced
      equation~\eqref{eq:red}.  We then split by $x \bmod 4$:
      \begin{itemize}
        \item $x \equiv 0 \pmod{4}$ (i.e.\ $x \equiv 4 \pmod{12}$):
              mod-$4$ obstruction of Lemma~\ref{lem:mod4a} (gives $3y^2 \equiv 2$).
        \item $x \equiv 3 \pmod{4}$ (i.e.\ $x \equiv 7 \pmod{12}$):
              mod-$4$ obstruction of Lemma~\ref{lem:mod4b} (gives $3Z^2 \equiv 1$).
        \item $x \equiv 1 \pmod{12}$:
              Lemma~\ref{lem:case1e} via mod-$8$ (even $n$) or Jacobi symbol
              (odd $n$).
        \item $x \equiv 10 \pmod{12}$:
              Lemma~\ref{lem:case1f} via mod-$48$ exhaustive check,
              Jacobi-symbol argument ($x \equiv 22 \pmod{48}$); the
              sub-case $x \equiv 34 \pmod{48}$ remains open.
      \end{itemize}
      These four sub-cases together cover all $x \equiv 1 \pmod{12}$,
      $4 \pmod{12}$, $7 \pmod{12}$, and $10 \pmod{12}$, which exhaust
      $x \equiv 1 \pmod{3}$.
  \end{description}

  Every residue class except $x \equiv 34 \pmod{48}$ is covered by a
  verified obstruction.  The proof of Theorem~\ref{thm:main} is therefore
  complete in five of the six residue classes modulo~$12$; the sub-case
  $x \equiv 34 \pmod{48}$ within $x \equiv 10 \pmod{12}$ remains open.
\end{proof}

% ─────────────────────────────────────────────────────────────────────────────
\section{Summary of Obstructions}

\begin{center}
\renewcommand{\arraystretch}{1.25}
\begin{tabular}{p{3.2cm}p{5.2cm}p{5.6cm}}
\toprule
Residue class of $x$ & Method & Key obstruction \\
\midrule
$x \equiv 0 \pmod{3}$
  & Direct mod-$3$ (Lemma \ref{lem:mod3x0})
  & $y^2 \equiv 2 \pmod{3}$ impossible \\
$x \equiv 2 \pmod{3}$
  & 3-adic descent + mod-$9$ (Lemma \ref{lem:mod3x2})
  & $x^3 \equiv 8 \not\equiv 2 \pmod{9}$ \\
$x \equiv 4 \pmod{12}$
  & Mod-$4$ on~\eqref{eq:red} (Lemma \ref{lem:mod4a})
  & $3y^2 \equiv 2 \pmod{4}$ impossible \\
$x \equiv 7 \pmod{12}$
  & Mod-$4$ on~\eqref{eq:red} (Lemma \ref{lem:mod4b})
  & $3Z^2 \equiv 1 \pmod{4}$ impossible \\
$x \equiv 1 \pmod{12}$, $n$ even
  & Mod-$8$ (Lemma \ref{lem:case1e}, Step~2)
  & LHS $\equiv 3$, RHS $\equiv 7 \pmod{8}$ \\
$x \equiv 1 \pmod{12}$, $n$ odd
  & Jacobi symbol (Lemma \ref{lem:case1e}, Step~3)
  & $2$ not a QR mod prime $p \equiv 3,5 \pmod{8}$ \\
$x \equiv 10,46 \pmod{48}$
  & Exhaustive mod-$48$ (Lemma \ref{lem:case1f})
  & No solution in $\mathbb{Z}/48\mathbb{Z}$ \\
$x \equiv 22 \pmod{48}$
  & Jacobi symbol (Lemma \ref{lem:case1f})
  & $2$ not a QR mod $p \equiv 3 \pmod{8}$ \\
$x \equiv 34 \pmod{48}$
  & \emph{Open} (Lemma \ref{lem:case1f})
  & Structure: $y=2k$, $z=3w$, $w$ odd, $k^2\equiv 1\pmod{9}$;
    no contradiction found. \\
\bottomrule
\end{tabular}
\end{center}

\begin{remark}
  The three corrections made from an earlier incorrect draft are:
  \begin{enumerate}
    \item \emph{Lemma statement (Error~1):} The preliminary reduction lemma
          correctly concludes that any solution must have $x \equiv 1
          \pmod{3}$ (the two other classes being eliminated), not
          $x \equiv 0 \pmod{3}$ as misstated in that draft.
    \item \emph{Mod-$3$ reduction (Error~2):} When reducing the equation
          modulo~$3$, the term $xz^2$ must not be silently dropped;
          the identity $(*)$ avoids this by placing all terms on one side.
    \item \emph{3-adic descent (Error~3):} The claim that $3^n \mid z$ for
          all $n$ (forcing $z = 0$) via an inductive mod-$9$ argument is
          invalid: for $n \geq 1$ the term $3^{2n}xw^2$ is automatically
          divisible by~$9$, so the inductive step gives no constraint on $w$.
          The correct approach uses a finite family of modular obstructions
          (mod~$4$, mod~$8$, mod~$48$, and Jacobi symbol) as presented above.
  \end{enumerate}
  A subsequent proof attempt for the open sub-case $x \equiv 34 \pmod{48}$
  used the substitution $x = 2p$ to derive a reduced equation
  $p(4p^2-12k^2-9w^2) = 1-2k^2$ (Lemma~\ref{lem:structure10}), but
  contained a sign error in the central identity: equation~(1) was stated
  as $m(m^2-3y^2-z^2) = y^2-2$ rather than the correct $2-y^2$.
  Upon correcting the sign, Steps~1--3 of that proof remain valid and are
  incorporated as Lemma~\ref{lem:structure10}; however, Step~4 (a mod-$9$
  argument) no longer yields a contradiction --- it gives only $k^2 \equiv 1
  \pmod{9}$, which is satisfiable.  The sub-case $x \equiv 34 \pmod{48}$
  therefore remains \textbf{open}.
\end{remark}

\end{document}
